\section{The \texttt{MCPL\_input} McXtrace Component}
Source-like component that reads photon state parameters from an mcpl-file.

\subsection*{Identification}
\begin{itemize}
  \item \textbf{Author:} Erik B Knudsen, Gregory S Tucker and Peter Willendrup
  \item \textbf{Origin:} DTU Physics / European Spallation Source ERIC
  \item \textbf{Date:} Aug 2016, consolidated with McStas MCPL\_input\_once Oct 2026
\end{itemize}

\subsection*{Description}
Source-like component that reads photon state parameters from a binary mcpl-file.

MCPL is short for Monte Carlo Particle List, and is a format for sharing events between e.g. MCNP(X), Geant4 and McXtrace.

The --ncount given on the commandline is overwritten by (\#photons in the file) x repeat\_count. With preload=1 (and always on GPU) only the photons within [Emin, Emax] are counted.

When used with MPI, the file is split between the MPI processes: each process reads (and preloads) only its own slice of about (\#events in the file) / (\#MPI processes), so that every event in the file is used once per repetition across the whole run. (Before Oct 2026, MCPL\_input instead let every process read the whole file and smeared all events on the non-master processes.)

repeat\_count \textgreater{} 1 replays each process' slice that many times, with the event weight divided accordingly. All replays after the first are smeared using E\_smear, pos\_smear and dir\_smear (the first pass is too if always\_smear is set), so repeat\_count \textgreater{} 1 MUST be combined with at least one *\_smear option; if all are 0, repeat\_count is reset to 1.

This is the same implementation as McStas' MCPL\_input, adapted to photons.

Example: MCPL\_input(filename=voutput,verbose=1,repeat\_count=1,E\_smear=0.1,pos\_smear=0.001,dir\_smear=0.01)

\subsection*{Input parameters}
Parameters in \textbf{boldface} are required; the others are optional.

\begin{longtable}{p{0.22\textwidth}p{0.12\textwidth}p{0.46\textwidth}p{0.14\textwidth}}
\toprule
\textbf{Name} & \textbf{Unit} & \textbf{Description} & \textbf{Default} \\
\midrule
\endhead
filename & str & Name of photon mcpl file to read & 0 \\
polarisationuse &   & If !=0 read polarisation vectors from file & 1 \\
verbose &   & Print debugging information for the first 10 photons read (on the master process) & 0 \\
Emin & keV & Lower energy bound. Particles found in the MCPL-file below the limit are skipped & 0 \\
Emax & keV & Upper energy bound. Particles found in the MCPL-file above the limit are skipped & FLT\_MAX \\
repeat\_count & 1 & Replay the contents of the MCPL file this number of times, dividing the event weights by the same number. MUST be combined with \_smear options! & 1 \\
E\_smear & 1 & When replaying events, make a flat MC choice within E\_smear*E around particle energy E & 0 \\
pos\_smear & m & When replaying events, make a flat MC choice of position within pos\_smear around particle starting position & 0 \\
dir\_smear & deg & When replaying events, make an MC choice of direction within dir\_smear around particle direction & 0 \\
preload &   & Load particles during INITIALIZE. On GPU preload is forced & 0 \\
always\_smear &   & If !=0, also smear the events on their first pass through the file & 0 \\
\bottomrule
\end{longtable}

\subsection*{Links}
\begin{itemize}
  \item Component source code found in file \texttt{MCPL\_input.comp}.
\end{itemize}
\IfFileExists{misc/MCPL_input_static.tex}{\input{misc/MCPL_input_static.tex}}{}